% ============================================================
%  AML Project 4 -- IEEE Report
%  CSCI-6767 Applied Machine Learning & Data Analytics
%  Spring 2026
%  Team: Gabil Gurbanov, Hamida Hagverdiyeva
% ============================================================
\documentclass[conference]{IEEEtran}

\usepackage[T1]{fontenc}
\usepackage[utf8]{inputenc}
\usepackage{cite}
\usepackage{amsmath,amssymb,amsfonts}
\usepackage{graphicx}
\usepackage{textcomp}
\usepackage{xcolor}
\usepackage{booktabs}
\usepackage{url}
\usepackage{hyperref}

% Path to uploaded figures (upload model_cache/figs/ contents to Overleaf)
\graphicspath{{figs/}}

\begin{document}

% --- TITLE ---------------------------------------------------------------
\title{Non-Linearity, Tree-Based Classification, and SVM:\\
       A Comparative Study on Regression and Intrusion Detection}

\author{
  \IEEEauthorblockN{Gabil Gurbanov}
  \IEEEauthorblockA{CSCI-6767, Spring 2026 \\
    \textit{Fairleigh Dickinson University}}
  \and
  \IEEEauthorblockN{Hamida Hagverdiyeva}
  \IEEEauthorblockA{CSCI-6767, Spring 2026 \\
    \textit{Fairleigh Dickinson University}}
}

\maketitle

% --- ABSTRACT ------------------------------------------------------------
\begin{abstract}
This paper applies and compares a spectrum of machine learning methods
across two distinct tasks: non-linear regression on the Seoul Bike Sharing
Demand dataset and binary classification on the KDD Cup 1999 network
intrusion dataset. For regression we evaluate polynomial regression, step
functions, B-splines, natural cubic splines, locally weighted scatterplot
smoothing (LOESS), and generalized additive models (GAMs) against a linear
baseline. For classification we evaluate bagging, random forests, gradient
boosting, and support vector machines (SVMs) with four kernel types. GAMs
achieve the best regression performance (RMSE~$=$~317.70, $R^2$~$=$~0.758),
while tree ensembles and RBF-kernel SVMs dominate classification
($F_1 \geq 0.999$). Our analysis highlights why model choice must align with
data structure, scale, and interpretability requirements.
\end{abstract}

\begin{IEEEkeywords}
non-linear regression, generalized additive models, splines, random forest,
gradient boosting, support vector machines, network intrusion detection,
bike sharing demand
\end{IEEEkeywords}

% --- 1. INTRODUCTION -----------------------------------------------------
\section{Introduction}\label{sec:intro}

Linear models remain a practical baseline for many supervised learning tasks,
yet they impose a rigid assumption: the response is a weighted sum of the
input features. When the true data-generating process is non-linear, or when
decision boundaries are complex and high-dimensional, linear models
systematically underfit and yield poor generalisation.

This project pursues two complementary goals. First, we demonstrate the
gains achievable by replacing a linear regression baseline with progressively
flexible non-linear alternatives on a real-world demand-forecasting problem.
Second, we compare ensemble tree methods and kernel SVMs for a large-scale
binary classification task in the network-security domain.

The two-dataset strategy is deliberate: a single dataset forced into both
tasks would compromise the integrity of at least one comparison. Using
domain-appropriate datasets lets each model family operate under conditions
where it can reasonably be expected to succeed or fail, producing cleaner
insights.

% --- 2. DATASETS ---------------------------------------------------------
\section{Datasets}\label{sec:data}

\subsection{Seoul Bike Sharing Demand (Regression)}

The Seoul Bike Sharing Demand dataset \cite{seouldataset} is drawn from the
UCI Machine Learning Repository (ID~560). It contains 8,760 hourly records
spanning one full year, with 14 features including air temperature, humidity,
wind speed, visibility, dew-point temperature, solar radiation, snowfall,
rainfall, and categorical indicators for season, holiday, and functioning day.
The target variable is the hourly count of rented bicycles
(\textit{Rented Bike Count}).

\textbf{Why it fits.} Hourly demand exhibits strong non-linear patterns:
a bimodal peak at the morning and evening rush hours, a non-monotone
temperature response (rentals peak near 20\textdegree{}C and fall in extreme
heat), and seasonal periodicity. These characteristics make the dataset
well-suited for evaluating splines, LOESS, and GAMs against linear baselines.

\textbf{Preprocessing.} The \texttt{Date} column is dropped. Categorical
columns (\texttt{Seasons}, \texttt{Holiday}, \texttt{Functioning Day}) are
one-hot encoded. The target is log-transformed to reduce right skew. Features
are split 80/20 into training and test sets with \texttt{random\_state=42}.

\subsection{KDD Cup 1999 -- 10\% Subset (Classification)}

The KDD Cup 1999 dataset \cite{kddcup99} is a widely used benchmark for
intrusion-detection system (IDS) research. The 10\% subset contains 494,021
connection records and 41 features comprising basic connection attributes
(duration, protocol type, service, flag), content features (number of failed
logins, root accesses), and traffic features (same-host and same-service
connection rates). The original multi-class target is collapsed to a binary
label: \textit{normal} (0) vs.\ \textit{attack} (1).

\textbf{Why it fits.} The large sample size highlights scalability
differences between tree ensembles and kernel SVMs. The rich, mixed feature
space (continuous + categorical) allows all model families to be fairly
compared.

\textbf{Preprocessing.} Categorical features (\texttt{protocol\_type},
\texttt{service}, \texttt{flag}) are label-encoded. An 80/20 stratified
train/test split preserves class proportions. Features are standardised
(zero mean, unit variance) for SVM. Because kernel SVM training complexity
is $\mathcal{O}(n^2)$--$\mathcal{O}(n^3)$ in the number of samples $n$, SVMs
are trained on a 15,000-sample stratified subset; tree ensembles use the
full training set.

\begin{figure}[t]
  \centering
  \includegraphics[width=\columnwidth]{figs/bike_eda1.png}
  \caption{Seoul Bike dataset: hourly demand distribution (left) and
    temperature vs.\ rented bike count scatter (right), illustrating
    non-linear structure in both the time and temperature axes.}
  \label{fig:bike_eda1}
\end{figure}

\begin{figure}[t]
  \centering
  \includegraphics[width=\columnwidth]{figs/bike_corr.png}
  \caption{Pearson correlation heatmap for the Seoul Bike dataset.
    Temperature and dew-point temperature are nearly collinear
    ($r \approx 0.91$).}
  \label{fig:bike_corr}
\end{figure}

% --- 3. METHODS ----------------------------------------------------------
\section{Methods}\label{sec:methods}

\subsection{Non-Linear Regression Models}

\subsubsection{Baseline -- Linear Regression}
Ordinary least squares (OLS) fits the model
\begin{equation}
  \hat{y} = \beta_0 + \sum_{j=1}^{p} \beta_j x_j,
  \label{eq:ols}
\end{equation}
where $\hat{y}$ is the predicted log-transformed count and $x_j$ are
the pre-processed features. This serves as the lower-bound baseline.

\subsubsection{Polynomial Regression}
Feature space is expanded by including all polynomial terms up to degree $d$:
\begin{equation}
  \hat{y} = \beta_0 + \beta_1 x + \beta_2 x^2 + \cdots + \beta_d x^d.
  \label{eq:poly}
\end{equation}
Ridge regularisation ($\ell_2$ penalty) is applied via \texttt{RidgeCV} to
control variance at higher degrees. The optimal degree is selected by
cross-validated RMSE over $d \in \{2, 3, 4, 5\}$.

\subsubsection{Step Functions}
The range of each continuous predictor is partitioned into $K$ bins using
equal-frequency quantile cuts, generating indicator dummy variables
$\mathbf{1}(x \in \mathcal{C}_k)$ for $k = 1, \ldots, K$. A ridge regression
is then fit on these expanded features. Step functions allow abrupt level
shifts between intervals but cannot capture smooth transitions within them.

\subsubsection{B-Splines}
Scikit-learn's \texttt{SplineTransformer} generates a basis of B-spline
functions of degree 3 (cubic) with $K$ interior knots placed at uniform
quantiles. The transformed features are passed to ridge regression. B-splines
guarantee $C^2$ continuity everywhere (smooth first and second derivatives),
avoiding the artificial jumps of step functions while remaining piecewise
polynomial.

\subsubsection{Natural Cubic Splines}
Natural cubic splines (NCS) are cubic splines that are additionally
constrained to be linear beyond the boundary knots, reducing edge-effect
artefacts. The design matrix is constructed using \texttt{patsy}'s
\texttt{cr()} formula syntax:
\begin{equation}
  \hat{y} = \sum_{k=1}^{K} \alpha_k N_k(x),
  \label{eq:ncs}
\end{equation}
where $N_k$ are the natural spline basis functions. Ridge regression is
applied to the full multivariate NCS design matrix.

\subsubsection{LOESS (Local Regression)}
LOESS (locally estimated scatterplot smoothing) fits weighted local
polynomial regressions at each target point $x_0$:
\begin{equation}
  \hat{f}(x_0) = \arg\min_{\boldsymbol{\beta}}
    \sum_{i=1}^{n} K\!\left(\frac{x_i - x_0}{h}\right)
    \bigl(y_i - \beta_0 - \beta_1 x_i\bigr)^2,
  \label{eq:loess}
\end{equation}
where $K(\cdot)$ is a tri-cube kernel and $h$ controls the bandwidth. We
apply LOESS to temperature alone (fraction $= 0.10$) as a univariate
interpretive smoother. Because it is not a multivariate parametric model,
its test-set metrics are not directly comparable to the other regressors;
they are included to illustrate the technique and its limitations.

\subsubsection{Generalized Additive Model (GAM)}
A GAM decomposes the response into a sum of smooth functions of individual
predictors:
\begin{equation}
  g\!\bigl(\mathbb{E}[y]\bigr)
    = \beta_0 + \sum_{j=1}^{p} f_j(x_j),
  \label{eq:gam}
\end{equation}
where $g$ is the link function (identity for Gaussian regression) and each
$f_j$ is an arbitrary smooth function estimated via penalised regression
splines. We use \texttt{pygam.LinearGAM} with a spline term $s(\cdot)$ for
each numeric predictor. The penalty controls the roughness of each $f_j$,
making GAMs inherently regularised and interpretable through partial
dependence plots.

\subsection{Tree-Based Classification Methods}

\subsubsection{Bagging}
Bootstrap aggregating (bagging) \cite{breiman1996bagging} trains $B$
decision trees on independently bootstrapped training sets and averages
(or majority-votes) their predictions:
\begin{equation}
  \hat{f}_{\text{bag}}(x) = \frac{1}{B} \sum_{b=1}^{B} T_b(x),
  \label{eq:bagging}
\end{equation}
where $T_b$ is the tree trained on the $b$-th bootstrap sample. Variance is
reduced without increasing bias. We use $B = 50$ decision-tree base learners.

\subsubsection{Random Forest}
Random forests \cite{breiman2001rf} augment bagging by also randomising the
feature subset considered at each split. At each node, only $m = \lfloor
\sqrt{p} \rfloor$ features (for classification) are eligible, decorrelating
trees and further reducing ensemble variance. We train $B = 100$ trees.

\subsubsection{Gradient Boosting}
Gradient boosting \cite{friedman2001gbm} constructs an additive model
sequentially, each new tree $T_b$ fitting the negative gradient (pseudo
residuals) of the loss from the current ensemble:
\begin{equation}
  F_b(x) = F_{b-1}(x) + \eta T_b(x),
  \label{eq:gbm}
\end{equation}
where $\eta$ is the learning rate (shrinkage). We use $B = 100$ stages with
$\eta = 0.1$ and maximum tree depth of 3.

\subsection{Support Vector Machines}

An SVM \cite{cortes1995svm} finds the maximum-margin hyperplane
\begin{equation}
  \min_{\mathbf{w},b} \frac{1}{2}\|\mathbf{w}\|^2 + C \sum_{i=1}^{n} \xi_i
  \quad \text{s.t.} \quad y_i(\mathbf{w}^\top \phi(x_i)+b) \geq 1 - \xi_i,
  \label{eq:svm}
\end{equation}
where $\phi$ is a kernel-induced feature map, $C$ is the regularisation
parameter, and $\xi_i \geq 0$ are slack variables. We evaluate four kernels:

\begin{itemize}
  \item \textbf{Linear}: $K(x,x') = x^\top x'$
  \item \textbf{RBF}: $K(x,x') = \exp(-\gamma\|x-x'\|^2)$
  \item \textbf{Polynomial}: $K(x,x') = (\gamma x^\top x' + r)^d$, $d=3$
  \item \textbf{Sigmoid}: $K(x,x') = \tanh(\gamma x^\top x' + r)$
\end{itemize}

All SVMs use $C = 1$ and default scikit-learn gamma (\texttt{scale}).

% --- 4. EXPERIMENTS ------------------------------------------------------
\section{Experiments and Results}\label{sec:experiments}

\subsection{Regression Evaluation Metrics}

We report Root Mean Squared Error (RMSE), Mean Absolute Error (MAE), and the
coefficient of determination $R^2$:
\begin{equation}
  \text{RMSE} = \sqrt{\frac{1}{n}\sum_{i=1}^{n}(y_i - \hat{y}_i)^2}, \quad
  \text{MAE}  = \frac{1}{n}\sum_{i=1}^{n}|y_i - \hat{y}_i|,
  \label{eq:reg_metrics}
\end{equation}
\begin{equation}
  R^2 = 1 - \frac{\sum(y_i-\hat{y}_i)^2}{\sum(y_i-\bar{y})^2}.
  \label{eq:r2}
\end{equation}
Lower RMSE/MAE and higher $R^2$ indicate better fit.

\subsection{Regression Results}

\begin{table}[t]
  \caption{Regression Model Comparison on Seoul Bike Test Set}
  \label{tab:regression}
  \centering
  \begin{tabular}{lccc}
    \toprule
    \textbf{Model} & \textbf{RMSE} & \textbf{MAE} & $\mathbf{R^2}$ \\
    \midrule
    GAM (LinearGAM)            & \textbf{317.70} & 235.18 & \textbf{0.7578} \\
    Polynomial (degree 3)      & 350.99 & \textbf{215.75} & 0.7043 \\
    B-Splines (SplineTransformer) & 351.50 & 260.98 & 0.7035 \\
    Natural Cubic Splines      & 384.94 & 285.28 & 0.6444 \\
    Step Functions             & 400.64 & 299.26 & 0.6147 \\
    Baseline (Linear)          & 440.78 & 330.39 & 0.5337 \\
    LOESS (temp., frac=0.10)   & 545.09 & 388.10 & 0.2869 \\
    \bottomrule
  \end{tabular}
\end{table}

Table~\ref{tab:regression} summarises results. All major non-linear models
outperform the linear baseline, confirming that linear regression is
insufficient for this task. GAM achieves the best RMSE and $R^2$, while
polynomial degree 3 achieves the lowest MAE. B-splines and polynomial
regression are nearly tied overall. The LOESS result is expected to be weaker
because it operates on a single feature; its primary utility is visualisation
of the temperature effect (Fig.~\ref{fig:loess}).

\begin{figure}[t]
  \centering
  \includegraphics[width=\columnwidth]{figs/poly_compare.png}
  \caption{Polynomial regression degree sweep: RMSE on the validation set
    as a function of polynomial degree. Degree 3 gives the best trade-off
    between bias and variance.}
  \label{fig:poly}
\end{figure}

\begin{figure}[t]
  \centering
  \includegraphics[width=\columnwidth]{figs/spline_1d.png}
  \caption{B-spline fit (blue) vs.\ LOESS trend (orange) on the temperature
    axis. Splines capture the smooth non-linear temperature response while
    remaining stable at the boundary.}
  \label{fig:loess}
\end{figure}

\begin{figure}[t]
  \centering
  \includegraphics[width=\columnwidth]{figs/gam_pdp.png}
  \caption{GAM partial dependence plots. Each panel shows the learned smooth
    effect $\hat{f}_j(x_j)$ with 95\% confidence intervals. Hour and
    temperature exhibit clear non-monotone patterns.}
  \label{fig:gam_pdp}
\end{figure}

\begin{figure}[t]
  \centering
  \includegraphics[width=\columnwidth]{figs/reg_compare.png}
  \caption{Final regression model comparison: RMSE (left) and $R^2$ (right)
    on the held-out test set. GAM leads on both metrics.}
  \label{fig:reg_compare}
\end{figure}

\begin{figure}[t]
  \centering
  \includegraphics[width=\columnwidth]{figs/reg_avp.png}
  \caption{Actual vs.\ predicted scatter plots for selected regression models.
    GAM predictions cluster most tightly around the identity line.}
  \label{fig:reg_avp}
\end{figure}

\subsection{Classification Evaluation Metrics}

For binary classification we report accuracy, precision, recall, $F_1$ score,
and Area Under the ROC Curve (AUC-ROC):
\begin{equation}
  F_1 = \frac{2 \cdot \text{Precision} \cdot \text{Recall}}
             {\text{Precision} + \text{Recall}},
  \label{eq:f1}
\end{equation}
\begin{equation}
  \text{AUC} = \int_0^1 \text{TPR}(t) d\text{FPR}(t).
  \label{eq:auc}
\end{equation}
Because the class distribution skews toward the attack class, $F_1$ and
AUC-ROC are the primary metrics.

\subsection{Classification Results}

\begin{table}[t]
  \caption{Classification Model Comparison on KDD Cup 1999 Test Set}
  \label{tab:classification}
  \centering
  \resizebox{\columnwidth}{!}{%
  \begin{tabular}{lccccc}
    \toprule
    \textbf{Model} & \textbf{Acc.} & \textbf{Prec.} & \textbf{Rec.} & $\mathbf{F_1}$ & \textbf{AUC} \\
    \midrule
    Bagging (DT)      & \textbf{0.9998} & \textbf{0.9999} & 0.9998 & \textbf{0.9999} & \textbf{1.0000} \\
    Random Forest     & \textbf{0.9998} & \textbf{0.9999} & 0.9998 & \textbf{0.9999} & \textbf{1.0000} \\
    Gradient Boosting & 0.9997 & 0.9998 & 0.9997 & 0.9998 & 0.9999 \\
    SVM -- RBF        & 0.9983 & 0.9992 & 0.9987 & 0.9989 & 0.9999 \\
    SVM -- Polynomial & 0.9980 & 0.9989 & 0.9987 & 0.9988 & 0.9978 \\
    SVM -- Linear     & 0.9971 & 0.9983 & 0.9981 & 0.9982 & 0.9977 \\
    SVM -- Sigmoid    & 0.9496 & 0.9682 & 0.9690 & 0.9686 & 0.8861 \\
    \bottomrule
  \end{tabular}}
\end{table}

Table~\ref{tab:classification} presents results. Tree ensembles and the RBF
kernel dominate. Gradient boosting trails the parallel ensembles slightly,
because its sequential construction limits the breadth of diversity it can
exploit on this highly separable problem. Among SVMs the kernel choice is
decisive: RBF and polynomial kernels learn compact, curved boundaries that
tightly enclose attack clusters, whereas the sigmoid kernel imposes an
S-shaped decision surface that clips incorrectly through portions of the
normal class, dropping accuracy by roughly 5 percentage points and AUC by
0.11 compared to RBF.

\begin{figure}[t]
  \centering
  \includegraphics[width=\columnwidth]{figs/clf_conf_all.png}
  \caption{Confusion matrices for all seven classification models. The
    sigmoid SVM matrix shows the highest off-diagonal counts.}
  \label{fig:conf}
\end{figure}

\begin{figure}[t]
  \centering
  \includegraphics[width=\columnwidth]{figs/roc_all.png}
  \caption{ROC curves for all classifiers. Sigmoid SVM is visually
    distinguishable from the rest, which cluster near the top-left corner.}
  \label{fig:roc}
\end{figure}

\begin{figure}[t]
  \centering
  \includegraphics[width=\columnwidth]{figs/gb_deviance.png}
  \caption{Gradient boosting: training and test deviance (log-loss) as a
    function of boosting stage. The test curve plateaus after $\approx 40$
    stages, indicating good generalisation without overfitting.}
  \label{fig:gb_deviance}
\end{figure}

\begin{figure}[t]
  \centering
  \includegraphics[width=\columnwidth]{figs/rf_importance.png}
  \caption{Random forest feature importances (mean decrease in impurity).
    Traffic-rate features dominate, indicating that connection volume patterns
    are the most discriminative signal for intrusion detection.}
  \label{fig:rf_importance}
\end{figure}

% --- 5. ERROR ANALYSIS ---------------------------------------------------
\section{Error Analysis}\label{sec:error}

\subsection{Regression}

Fig.~\ref{fig:reg_resid} shows residual plots. The linear baseline exhibits
a systematic fan-shaped pattern, confirming heteroscedastic errors -- a clear
sign of model misspecification from unmodelled non-linearity. GAM residuals
are notably more homoscedastic, though moderate underestimation persists at
peak demand hours (rush hours), where the distribution of counts is multi-modal
within each time bin and no single smooth function captures the full spread.

The near-perfect correlation between temperature and dew-point temperature
($r = 0.91$, Fig.~\ref{fig:bike_corr}) introduces mild multicollinearity in
polynomial and spline models; GAM's per-feature smooth functions sidestep this
issue by modelling each predictor independently.

LOESS underperforms because the univariate smoother cannot account for
confounding variables. When evaluated on the full test set, its predictions
are based solely on temperature and ignore hour-of-day, season, and weather
features -- making the reported RMSE an indicator of single-feature predictive
power rather than model quality per se.

\begin{figure}[t]
  \centering
  \includegraphics[width=\columnwidth]{figs/reg_resid.png}
  \caption{Residual plots (predicted vs.\ residual) for the linear baseline
    (left) and GAM (right). GAM residuals are more centred and homoscedastic.}
  \label{fig:reg_resid}
\end{figure}

\subsection{Classification}

The near-perfect tree ensemble metrics reflect a known characteristic of
the KDD Cup 1999 dataset: it contains a large number of duplicate or
near-duplicate attack records derived from synthetic traffic replay. As a
result, the effective diversity of attack patterns is lower than the raw
count suggests, making the problem easier than real-world IDS scenarios.
We acknowledge this limitation and note that evaluation on a modern benchmark
(e.g.\ UNSW-NB15 \cite{unswnb15}) would be necessary to validate
generalisation.

The sigmoid kernel's poor performance is attributable to its functional form.
Unlike the RBF kernel, which can represent arbitrary closed decision regions,
the sigmoid kernel is conditionally positive-definite only for specific
parameter settings. When $\gamma$ and $r$ are at default values, the resulting
kernel matrix may not be strictly positive-definite, leading to a poorly
conditioned optimisation and suboptimal boundaries.

The SVM subset (15,000 samples) introduces a potential bias: SVMs may
benefit from seeing more examples of rare attack subtypes. Tree ensembles,
which were trained on the full set ($\sim\!395,000$ training examples), have
no such constraint, making a direct performance comparison slightly unfair.
Nonetheless, the scale of the gap ($F_1 > 0.999$ for ensembles vs.\
$F_1 = 0.9989$ for RBF-SVM) is large enough to be meaningful.

% --- 6. CONCLUSION -------------------------------------------------------
\section{Conclusion}\label{sec:conclusion}

We evaluated nine models across two real-world tasks and draw the following
conclusions:

\begin{itemize}
  \item \textbf{Non-linear regression beats linear modelling.} Every major
    non-linear method outperforms OLS on Seoul Bike demand, with GAM achieving
    the best balance of accuracy and interpretability ($R^2 = 0.758$ vs.\
    $R^2 = 0.534$ for the linear baseline).
  \item \textbf{GAMs offer the best of both worlds.} By decomposing the
    response into interpretable per-feature curves (Fig.~\ref{fig:gam_pdp}),
    GAMs combine predictive power comparable to splines with the transparency
    needed for operational deployment.
  \item \textbf{Tree ensembles are highly effective for structured tabular
    classification.} Bagging and random forest both achieve $F_1 \geq 0.9999$
    while training on the full dataset and remaining scalable.
  \item \textbf{Kernel selection is critical for SVMs.} The gap between RBF
    ($F_1 = 0.9989$, AUC~$= 0.9999$) and sigmoid
    ($F_1 = 0.9686$, AUC~$= 0.8861$) underscores that the kernel must match
    the geometry of the problem.
  \item \textbf{For deployment:} we recommend GAM for interpretable regression
    and Random Forest for large-scale classification, where scalable training,
    deterministic feature importance, and high accuracy are all required.
\end{itemize}

% --- 7. TEAM CONTRIBUTIONS -----------------------------------------------
\section{Team Contributions}\label{sec:contributions}

\begin{itemize}
  \item \textbf{Gabil Gurbanov}: Task 1 -- non-linear regression model
    implementation (polynomial, step functions, B-splines, NCS, LOESS, GAM),
    regression visualisations and diagnostics, model comparison and analysis,
    Streamlit UI architecture and caching system.
  \item \textbf{Hamida Hagverdiyeva}: Tasks 2 \& 3 -- classification model
    implementation (bagging, random forest, gradient boosting, four SVM
    kernels), ROC curve and confusion matrix analysis, feature importance
    analysis, classification result interpretation.
\end{itemize}

% --- REFERENCES ----------------------------------------------------------
\begin{thebibliography}{9}

\bibitem{seouldataset}
S.\ Sathishkumar V, J.\ Park, and Y.\ Cho,
``Predicting Residential Energy Consumption Using CNN-LSTM Neural Networks,''
\textit{UCI Machine Learning Repository}, Dataset ID~560, 2020.
[Online]. Available: \url{https://archive.ics.uci.edu/ml/datasets/Seoul+Bike+Sharing+Demand}

\bibitem{kddcup99}
DARPA, ``KDD Cup 1999 Data,''
\textit{UCI Machine Learning Repository}, 1999.
[Online]. Available: \url{http://kdd.ics.uci.edu/databases/kddcup99/kddcup99.html}

\bibitem{breiman1996bagging}
L.\ Breiman, ``Bagging Predictors,''
\textit{Machine Learning}, vol.\ 24, no.\ 2, pp.\ 123--140, 1996.

\bibitem{breiman2001rf}
L.\ Breiman, ``Random Forests,''
\textit{Machine Learning}, vol.\ 45, no.\ 1, pp.\ 5--32, 2001.

\bibitem{friedman2001gbm}
J.\ H.\ Friedman, ``Greedy Function Approximation: A Gradient Boosting Machine,''
\textit{Annals of Statistics}, vol.\ 29, no.\ 5, pp.\ 1189--1232, 2001.

\bibitem{cortes1995svm}
C.\ Cortes and V.\ Vapnik, ``Support-Vector Networks,''
\textit{Machine Learning}, vol.\ 20, no.\ 3, pp.\ 273--297, 1995.

\bibitem{unswnb15}
N.\ Moustafa and J.\ Slay, ``UNSW-NB15: A Comprehensive Data Set for Network
Intrusion Detection Systems,'' in \textit{Proc.\ Military Communications and
Information Systems Conf.\ (MilCIS)}, Canberra, Australia, 2015.

\bibitem{hastie2009esl}
T.\ Hastie, R.\ Tibshirani, and J.\ Friedman,
\textit{The Elements of Statistical Learning}, 2nd ed.\
New York: Springer, 2009.

\bibitem{james2021isl}
G.\ James, D.\ Witten, T.\ Hastie, and R.\ Tibshirani,
\textit{An Introduction to Statistical Learning}, 2nd ed.\
New York: Springer, 2021.

\end{thebibliography}

\end{document}
